We evaluate our proposed methods using two variants of \texttt{Qwen2.5-Instruct}~(3B and 7B) in Table~\ref{T2}. 
HyperGraphPro consistently achieves the best average performance, showing remarkable improvements over both naive generation and retrieval-augmented baselines.
With \texttt{Qwen2.5-3B-Instruct}, our method reaches an average F1 of 54.56, compared to 49.89 from the strongest baseline (Graph-R1).
Similarly, under the \texttt{Qwen2.5-7B-Instruct}, our method achieves the highest overall performance with an average F1 of 59.38, improving upon the best baseline (Graph-R1, 55.94) by more than 3.4 points.
Note that the performance improvement of HyperGraphPro is relatively higher on multi-hop question answering datasets such as 2Wiki., HotpotQA, and MuSiQue, compared to NQ dataset.
These results highlight that our step-aware graph evolution and policy optimization helps more effective reasoning on complex question answering tasks.


\subsection{Additional Experimental Results}
\label{sec:additional}
\begin{table}[t]
% \begin{minipage}[t!]{0.49\textwidth}
    % \setlength{\tabcolsep}{0pt} 
    % \small
    % \begin{tabular}{l|C{1.3cm}C{1.3cm}C{1.3cm}C{1.3cm}|C{1.3cm}C{1.3cm}C{1.3cm}|C{1.3cm}}
    \begin{center}
    \resizebox{\columnwidth}{!}{
    \begin{tabular}{c|cccccc|cccccc}
        \toprule
           & \multicolumn{2}{c}{2Wiki.}&\multicolumn{2}{c}{HotpotQA}&\multicolumn{2}{c|}{MuSiQue}& \multicolumn{2}{c}{2Wiki.}&\multicolumn{2}{c}{HotpotQA}&\multicolumn{2}{c}{MuSiQue} \\
           \cmidrule(l{3pt}r{3pt}){2-3} \cmidrule(l{3pt}r{3pt}){4-5} \cmidrule(l{3pt}r{3pt}){6-7}\cmidrule(l{3pt}r{3pt}){8-9} \cmidrule(l{3pt}r{3pt}){10-11} \cmidrule(l{3pt}r{3pt}){12-13}
         Method   & EM & F1 &EM & F1 & EM & F1 & EM & F1 &EM & F1 & EM & F1  \\
            \midrule
        &\multicolumn{6}{l|}{\textbf{\textit{Qwen3-0.6B}}} & \multicolumn{6}{l}{\textbf{\textit{Qwen3-1.7B}}} \\
           Search-R1  & 22.66 &29.28  & 31.25& 40.99   & 7.81 & 12.82 & 31.25 & 36.79  & 39.06& 44.58  & 10.94 & 17.23\\
           Graph-R1   &{24.22} &36.03  & 32.81 & 42.21 & 19.53 & 27.38 &{42.97} & 46.55  & 46.88 & 50.25 &29.69 & 37.72\\
            \rowcolor{HyperGraphPro!15}
          \textbf{HyperGraphPro}  
          &\textbf{28.91} &\textbf{40.72} & \textbf{36.72} &\textbf{47.39} & \textbf{23.44} & \textbf{30.42} &\textbf{47.66} &\textbf{50.91} & \textbf{50.78} & \textbf{55.02} & \textbf{32.81} & \textbf{40.73}\\
          
        \bottomrule
    \end{tabular}
    }
    \end{center}
    \vspace{-10pt}
    \caption{Experimental results using Qwen3-0.6B     and Qwen3-1.7B.
    }
    \label{tab:qwen3}
    \vspace{-10pt}
\end{table}

\begin{table}[t]
\begin{minipage}[t!]{0.49\textwidth}
    \begin{center}
    \resizebox{\columnwidth}{!}{
    \begin{tabular}{cccccccc}
        \toprule
           & & \multicolumn{2}{c}{2Wiki.}&\multicolumn{2}{c}{MuSiQue} \\
           \cmidrule(l{3pt}r{3pt}){3-4} \cmidrule(l{3pt}r{3pt}){5-6}
            Str. guided ret. &  Progress-aware PO & EM & F1 &EM & F1 \\
            \midrule
           &   & 50.00 &57.56  & 32.81& 40.51 \\
           \checkmark &   &{51.56} &58.45  & 33.59 & 42.14\\
            & \checkmark  &{53.91} &60.47  & 34.38& 45.16 \\
            \rowcolor{HyperGraphPro!15}
          \checkmark & \checkmark  
          &\textbf{55.47} &\textbf{61.72} & \textbf{37.50} &\textbf{47.27} \\
          
        \bottomrule
    \end{tabular}
    }
    \end{center}
    \vspace{-10pt}
    \caption{Ablation studies of our HyperGraphPro on multi-hop question answering datasets with Qwen-2.5-3B. (Str. guided ret.: Structure-guided hypergraph retrieval, Progress-aware PO: Progress-aware Step-level policy optimization).
    }
    \label{tab:abl2}
    \vspace{-10pt}
    % \vspace{-3mm}
    % }
\end{minipage}
\hfill
\begin{minipage}[t!]{0.50\textwidth}
    \begin{center}
    \resizebox{\columnwidth}{!}{
    \begin{tabular}{cccccccc}
        \toprule
           & & \multicolumn{2}{c}{2Wiki.} &\multicolumn{2}{c}{MuSiQue} \\
           \cmidrule(l{3pt}r{3pt}){3-4} \cmidrule(l{3pt}r{3pt}){5-6} 
          Step Progress & Structure Progress    & EM & F1    & EM & F1  \\
         \midrule 
         % \midrule
              &  & 51.56 & 58.45  & 33.59 & 42.14 \\
         \checkmark &      & 53.13 & 59.61 & 35.16 & 44.81 \\
          & \checkmark      & 54.69 & 60.55  & 34.38 & 44.29 \\
         \midrule
         \rowcolor{HyperGraphPro!15}
         \checkmark & \checkmark  &\textbf{55.47} &\textbf{61.72} & \textbf{37.50} &\textbf{47.27} \\
        \bottomrule
    \end{tabular}
    }
    \end{center}
    \vspace{-10pt}
    \caption{Performance comparison according to the policy optimization using Qwen-2.5-3B. Step progress means step progress-based dense rewarding. Structure progress means structure-consistent progressive dense rewarding.
    }
    \label{tab:po}
    \vspace{-10pt}
\end{minipage}
\vspace{-10pt}
% \vspace{-20pt}
\end{table}


\noindent\textbf{Experimental results on Qwen3.}
To assess the effectiveness of our HyperGraphPro on other variants of LLMs, we conduct the experiments using \texttt{Qwen3-0.6B} and \texttt{Qwen3-1.7B}~(Table~\ref{tab:qwen3}).
From the table, our HyperGraph outperforms other baseline methods using Qwen3 models.
This experimental result demonstrates that our proposed method is robust to small size variants of large language models.

\noindent\textbf{Ablation studies.}
To evaluate the effectiveness of the structure-guided hypergraph retrieval~(Sec.~\ref{sec:ret}) and progress-aware policy optimization~(Sec.~\ref{sec:po}), we conduct an ablation study on two multi-hop QA benchmarks in Table~\ref{tab:abl2}.
The results show that both components contribute to performance improvements across all datasets.
Specifically, applying progress-aware stepwise policy optimization leads to better learning dynamics and enhanced reasoning ability.
Similarly, incorporating our structure-guided hypergraph retrieval framework alone yields consistent gains over the baseline, demonstrating the benefit of integrating hypergraph structural information during retrieval.
Then, combining both components, our full model achieves the best performance, with notable improvements in F1 scores, such as +4.16 on 2Wiki. and +6.76 on MuSiQue over the baseline.
These results confirm that both finer-grained policy updates and structure-aware retrieval are crucial for improving multi-hop question answering with large language models.

\begin{table}[t]
\begin{minipage}[t!]{0.49\textwidth}
    \begin{center}
    \resizebox{\columnwidth}{!}{
    \begin{tabular}{lcccccc}
        \toprule
           & \multicolumn{2}{c}{2Wiki.}&\multicolumn{2}{c}{MuSiQue} \\
           \cmidrule(l{3pt}r{3pt}){2-3} \cmidrule(l{3pt}r{3pt}){4-5} 
          Retrieval method    & EM & F1     & EM & F1  \\
         \midrule 
         % \midrule
         No Retrieval      & 20.31& 28.45 &  3.12 & 8.07 \\
         Knowledge corpus~(Ctxt.) & 31.25 & 43.84 & 3.91 & 7.65  \\
         Hyperedge~(Ctxt.) & 50.78 & 58.13  & 32.03 & 40.31 \\
         Graph-R1~(Ctxt.)  & 53.91 & 61.21  &  34.38 & 43.41 \\
         \rowcolor{HyperGraphPro!15}
         \midrule
         \textbf{HyperGraphPro~(\textit{ours})~(Ctxt.+Struct.)} &\textbf{55.47} &\textbf{61.72} & \textbf{37.50} &\textbf{47.27} \\
        \bottomrule
    \end{tabular}
    }
    \end{center}
    \vspace{-10pt}
    \caption{Performance comparison according to the retrieval methods on multi-hop question answering datasets using Qwen-2.5-3B. Ctxt. indicates the natural language context-based relevance. Struct. indicates the graph structure-based relevance.
    }
    \label{tab:ret}
    \vspace{-10pt}
\end{minipage}
\hfill
\begin{minipage}[t!]{0.50\textwidth}
    \begin{center}
    \resizebox{\columnwidth}{!}{
    \begin{tabular}{lcccccc}
        \toprule
           & \multicolumn{2}{c}{2Wiki.}&\multicolumn{2}{c}{MuSiQue} \\
           \cmidrule(l{3pt}r{3pt}){2-3} \cmidrule(l{3pt}r{3pt}){4-5} 
          Method    & \# Turns & F1     & \# Turns & F1  \\
         \midrule 
         % \midrule
         \textbf{\textit{Qwen2.5-3B-Instruct}}\\
         Graph-R1  & 2.63 & 57.56  &  3.79 & 40.51 \\
         \rowcolor{HyperGraphPro!15}
         \textbf{HyperGraphPro~(\textit{ours})} &\textbf{2.45} &\textbf{61.72} & \textbf{3.16} &\textbf{47.27} \\
         \midrule
         \textbf{\textit{Qwen2.5-7B-Instruct}}\\
                  Graph-R1  & 2.72 & 65.04  &  3.92 & 46.17 \\
         \rowcolor{HyperGraphPro!15}
         \textbf{HyperGraphPro~(\textit{ours})} &\textbf{2.56} &\textbf{69.75} & \textbf{3.54} &\textbf{49.47} \\
        \bottomrule
    \end{tabular}
    }
    \vspace{-10pt}
    \end{center}
    \caption{Comparison on the number of turns between Graph-R1 and HyperGraphPro~(Ours).
    }
    \label{tab:eff}
    \vspace{-10pt}
\end{minipage}
\end{table}


\begin{table}[!t]
\tiny
    \begin{center}
    \resizebox{\textwidth}{!}{
    \begin{tabular}{p{2.8cm}|p{6.2cm}|p{5.4cm}}
        \toprule
        \textbf{Question \& Answer} &  \textbf{Graph-R1} & \textbf{HyperGraphPro~(Ours)} \\
        \midrule
          % \begin{tabular}{@{}l@{}}
          \textbf{Question: }When is election day for senate in the state where Taylor Hicks is from?
          \newline \textbf{Answer: }December 12, 2017
            % \end{tabular}
            &
\textcolor{blue}{$\langle$think$\rangle$}
To answer this question, I need to determine the state Taylor Hicks is from and then find out when the election day for the senate is in that state.
\textcolor{blue}{$\langle$/think$\rangle$}\newline
\textcolor{orange}{$\langle$query$\rangle$}
where is Taylor Hicks from
\textcolor{orange}{$\langle$/query$\rangle$}\newline
\textcolor{green}{$\langle$knowledge$\rangle$}
According to season five winner Taylor Hicks, who is from the state of Alabama, ...
\textcolor{green}{$\langle$/knowledge$\rangle$}\newline
\textcolor{blue}{$\langle$think$\rangle$}
The knowledge base states that Taylor Hicks, the season five winner, is from Alabama. Now I need to find out when the election day for the senate is in Alabama.
\textcolor{blue}{$\langle$/think$\rangle$}\newline
\textcolor{orange}{$\langle$query$\rangle$}
election day for senate in Alabama
\textcolor{orange}{$\langle$/query$\rangle$}\newline
\textcolor{green}{$\langle$knowledge$\rangle$}
A special election for the United States Senate in Alabama took place on December 12, 2017,...
\textcolor{green}{$\langle$/knowledge$\rangle$}\newline
\textcolor{blue}{$\langle$think$\rangle$}
The knowledge base provides information about a special election for the United States Senate in Alabama... Typically, senate elections occur every two years on the same day as the general midterm elections, which is usually the first Tuesday after the first Monday in November.
\textcolor{blue}{$\langle$/think$\rangle$}
\newline\textcolor{orange}{$\langle$query$\rangle$}
when is the election day for senate in the united states
\textcolor{orange}{$\langle$/query$\rangle$}\newline
\textcolor{green}{$\langle$knowledge$\rangle$}
A special election for the United States Senate in Alabama took place on December 12, 2017,...
The 2018 United States House of Representatives elections will be held on November 6, 2018.
\textcolor{green}{$\langle$/knowledge$\rangle$}\newline
\textcolor{blue}{$\langle$think$\rangle$}
Based on the information provided, the typical election day for the senate in the United States is the first Tuesday after the first Monday in November every two years. This is the general midterm election day.
\textcolor{blue}{$\langle$/think$\rangle$}\newline
\textcolor{red}{$\langle$answer$\rangle$}
The election day for the senate in the United States is the first Tuesday after the first Monday in November.
\textcolor{red}{$\langle$/answer$\rangle$}
            &
\textcolor{blue}{$\langle$think$\rangle$}
To answer this question, I need to find out which state Taylor Hicks is from first. Then, I can find the election day for the senate election in that state.
\textcolor{blue}{$\langle$/think$\rangle$}\newline
\textcolor{orange}{$\langle$query$\rangle$}
Taylor Hicks state
\textcolor{orange}{$\langle$/query$\rangle$}\newline
\textcolor{green}{$\langle$knowledge$\rangle$}
According to season five winner Taylor Hicks, who is from the state of Alabama, ...
\textcolor{green}{$\langle$/knowledge$\rangle$}\newline
\textcolor{blue}{$\langle$think$\rangle$}
From the search results, we can see that Taylor is from the state of Alabama. Now, I need to find out the election day for the senate election in Alabama.
\textcolor{blue}{$\langle$/think$\rangle$}\newline
\textcolor{orange}{$\langle$query$\rangle$}
election day for senate election in Alabama
\textcolor{orange}{$\langle$/query$\rangle$}\newline
\textcolor{green}{$\langle$knowledge$\rangle$}
A special election for the United States Senate in Alabama took place on December 12, 2017,...
\textcolor{green}{$\langle$/knowledge$\rangle$}\newline
\textcolor{blue}{$\langle$think$\rangle$}
From the search results, the special election for the United States Senate in Alabama took place on December 12, 2017.
\textcolor{blue}{$\langle$/think$\rangle$}\newline
\textcolor{red}{$\langle$answer$\rangle$}
December 12, 2017
\textcolor{red}{$\langle$/answer$\rangle$}\\
        \bottomrule
    \end{tabular}
    }
    \end{center}
    \vspace{-10pt}
    \caption{Comparison on outputs generated by Graph-R1 and HyperGraphPro~(Ours).
    }
    \label{tab:qual}
    \vspace{-12pt}
    % \vspace{-5mm}
    % }
\end{table}

%          
% Answer: 



\noindent\textbf{Policy optimization analysis.}
Additionally, we conduct an ablation study to examine the effect of different progress-aware dense rewarding functions used for policy optimization.
Table~\ref{tab:po} reports the experimental results on 2Wiki and MuSiQue.
The best performance is obtained when both step progress-based $r_t^{sp}$ and structure-consistent progressive $r_t^{struct}$ dense rewarding are jointly applied.
% On 2WikiMultiHopQA, this setting improves EM from 51.56 to \textbf{55.47} and F1 from 58.45 to \textbf{61.72}.
% Similarly, on MuSiQue, EM and F1 increase to \textbf{37.50} and \textbf{47.27}, respectively. 
These results suggest that the two dense reward signals capture complementary aspects of reasoning progress, and that their combination more effectively guides the policy toward accurate multi-hop reasoning.

\noindent\textbf{Retrieval method analysis.}
Table~\ref{tab:ret} compares different retrieval strategies on multi-hop QA datasets using \texttt{Qwen2.5-3B-Instruct}.
Our HyperGraphPro consistently achieves the best performance across all datasets, highlighting the effectiveness of integrating both contextual and structural relevance in retrieval.
Compared to the text-only retrieval baselines (Knowledge corpus and Hyperedge), HyperGraphPro yields substantial gains, demonstrating that leveraging graph structures enables more accurate and semantically coherent knowledge selection.
These results confirm that combining semantic and structural information is crucial for improving reasoning consistency in complex question answering.

\noindent\textbf{Efficiency analysis.}
We evaluate the efficiency of HyperGraphPro by comparing the average number of reasoning turns with Graph-R1 using \texttt{Qwen-2.5-3B-Instruct} and \texttt{Qwen-2.5-7B-Instruct}  in Table~\ref{tab:eff}.
The results show that HyperGraphPro achieves higher performance with fewer turns even though turn count is not used as a reward signal.
This demonstrates that our method improves both reasoning effectiveness and efficiency.

\subsection{Qualitative Analysis}
Here, we conduct a qualitative comparison between Graph-R1 and HyperGraphPro to illustrate how our reasoning framework enhances multi-hop reasoning (Table~\ref{tab:qual}).
While Graph-R1 correctly retrieves relevant knowledge snippets, it often produces redundant or conflicting reasoning steps—such as repeatedly querying the election day in Alabama and overgeneralizing the final answer to the national level.
In contrast, HyperGraphPro successfully generates the reasoning trajectory by progressively refining each step, leveraging retrieved graph context to eliminate unnecessary or misleading hops.
As shown in the example, HyperGraphPro accurately identifies the specific “special election” event on December 12, 2017, by integrating contextual knowledge about the Senate vacancy, while Graph-R1 misleads itself toward the general midterm schedule.
This highlights that HyperGraphPro effectively mitigates reasoning drift and enables precise, context-aware multi-hop inference.